\documentclass[四库全书文渊阁本]{ltc-guji}
\usepackage{comment}
\无标点模式
\setmainfont{TW-Kai}


\title{欽定四庫全書簡明目錄}

\begin{document}
\begin{封面}[底色={231, 198, 83}]
  \四库全书书名签{經部\空格[4] 簡明目錄卷一}
\end{封面}

\begin{空白页}
  \文本框[
    floating=true,
    x=5cm,
    y=10cm,
    font-size=28pt,
    grid-height=28pt,
    height=8]{檢討\臣 翁樹培覆勘}
  \文本框[
    floating=true,
    x=13cm,
    y=15cm,
    font-size=28pt,
    grid-height=28pt,
    height=11]{協勘官候補洗馬\臣 劉權之}
  \文本框[
    floating=true,
    x=14.5cm,
    y=16.8cm,
    font-size=28pt,
    grid-height=28pt,
    height=10]{\样式[grid-height=36pt]{謄錄舉人}\臣 李學沅}
\end{空白页}

\begin{正文}
\chapter{聖諭}
欽定四庫全書簡明目錄

聖諭

\begin{段落}[indent=2]
乾隆三十九年七月二十五日大學士于敏中等奉
\end{段落}

\begin{段落}[indent=1, first-indent=0]
諭旨四庫全書處總目於經史子集內分晰應刻應抄%
及應存書名三項各條下俱經撰有提要將一書原%
委撮舉大凡并詳著書人世次爵里可以一覽了然%
較之崇文總目範羅既廣體例加詳自應如此辦理%
第此次各省搜訪書籍有多至百種以上至六七百%
種者如浙江范懋柱等家其裒集收藏深可嘉尚前%
已降旨分別頒賞古今圖書集成及初印佩文韻府%
并擇其書尤雅者製詩親題卷端俾其子孫世守以%
為稽古藏書者勸今進到之書於攥輯後仍須發還%
本家而所撰總目若不載明係何人所藏則閱者不%
能知其書所自來亦無以彰各家珍弃資益之善著%
通查各省進到之書其一人而收藏百種以上者可%
稱為藏書之家即應將其姓名附載于各書提要末%
其在百種以下者亦應將由其省督撫主人採訪所%
得附載於後其官板刊刻及各處陳設庫貯者俱載%
內府所藏使其眉目分明更為詳細至現辦四庫全%
書總目提要多至萬餘種卷帙甚繁將來抄刻成書%
總閱已頗為不易自應於提要之外別刊簡明書目%
一編祗載其書若干卷註其朝代人撰則篇幅不繁%
而檢查較易俾學者由書目而尋提要由提要而得%
全書嘉與海內之士考訂源流用昭我朝文治之盛%
著四庫全書處總裁等遵照悉心妥辦並著通諭知%
之欽此
\end{段落}

\end{正文}
\end{document}
